\section{Hypothesis subversion: inflow paradigm and cosmological extensions}
\label{sec:inflow-hypotheses}

\begin{remark}[Status of this section]
The following propositions are hypothesis-grade. They originate in an
internal draft version (the ``inflow paradigm'') and have not yet been
mathematically proved. They are retained here as structured ideas for
future development.
\end{remark}

\subsection{Inflow picture as cascade dynamics}

\begin{proposition}[Inflow interpretation of the cascade, hypothesis-grade]
\label{prop:inflow-cascade}
\emph{(Hypothesis-grade.)}
The draft's ``inflow paradigm'' is structurally identical with the
cascade dynamics $S^8\to S^4\to S^2$:
\begin{itemize}
  \item ``information inflow'' = cascade flow from $S^8$ to $S^2$
  \item ``resistance as mass'' = the fraction $\bsym^2 \approx 0.302$
    that runs through $S^4$ (matter stratum) instead of directly to $S^2$
  \item ``biquaternionic core'' = the biquaternionic structure
    $q_n = it\cdot\mathbf{1} + z_i e_i \in \HH\otimes\mathbb{C}$
\end{itemize}

No new mathematical foundation is required; the inflow picture is an
intuitive reformulation of the existing cascade probabilities.
\end{proposition}

\paragraph{Transition from inflow to source-return closure.}
The inflow paradigm alone does not yet determine how the terminal OP sector should be read. In the present framework, OP~5 is not interpreted as an unrelated endpoint observable, but as the return-side readout of the same structural route that begins at OP~1. For that reason, the primary OP question is not raw endpoint coincidence, but whether the two-pass source-return transport remains inside one and the same native closure-bearing carrier class. This leads naturally to the following closure statement for the OP~1--OP~5 pair.


\subsection{OP transport closure, branch admissibility, and OP3 shell admissibility}
\label{sec:compressed-op-branch-shell}

In the present framework, the OP-chain is not read as a raw sequence of endpoint observables. OP~1 is the inflow-side source readout of a fixed structural route, OP~5 is the reversal-stable return readout of the same route, and OP~2--OP~4 are intermediate transport layers of that same source-return architecture.

Let $\Gamma$ be a fixed structural route in the admissible closure sector $\mathcal{A}_{\mathrm{cl}}\subseteq\mathcal{C}$, let
\[
T_{\Gamma}=A_{4}\circ A_{3}\circ A_{2},
\]
and let
\[
\rho_{\theta}:\mathcal{C}\to\mathcal{W}
\]
be the frame-dependent readout family. Let $K_{\mathrm{ncl}}$ denote the native closure carrier class, and let $\mathcal{I}$ be a closure-compatible observable invariant.

\begin{theorem}[Compressed OP branch-shell theorem]
\label{thm:compressed-op-branch-shell}
Assume that the route is HC-confined, two-pass class consistent, and intermediately class-admissible, and that $\mathcal{I}$ factors through the native closure carrier class and is linear on the admissible return and routed sectors. Then the following hold.
\begin{enumerate}[leftmargin=2em,label=(\roman*)]
\item \textbf{OP closure.}
\[
\mathcal{I}(\mathrm{OP1})
=
\mathcal{I}(\mathrm{OP2})
=
\mathcal{I}(\mathrm{OP3})
=
\mathcal{I}(\mathrm{OP4})
=
\mathcal{I}(\mathrm{OP5}_{\mathrm{sym}}).
\]
\item \textbf{Admissible branch preservation.} For every closure-admissible branch $b$,
\[
\mathcal{I}(\mathrm{OP5}^{(b)}_{\mathrm{sym}})=\mathcal{I}(\mathrm{OP1}).
\]
\item \textbf{OP3 shell-admissible preservation.} For every OP~3 shell-admissible refinement $(r,s)$,
\[
\mathcal{I}(\rho_{\theta}(c_{3}^{(r,s)}))=\mathcal{I}(\mathrm{OP1}).
\]
\item \textbf{Combined routed preservation.} If admissible branches and shell-admissible refinements are routed with normalized admissible weights, then the resulting visible return still satisfies
\[
\mathcal{I}(\mathrm{OP5}_{\mathrm{bs\text{-}rout}})=\mathcal{I}(\mathrm{OP1}).
\]
\end{enumerate}
\end{theorem}

\begin{proof}
Statements (i)--(ii) are the compressed consequences of Theorem~\ref{thm:op15-closure}, Theorem~\ref{thm:admissible-branch}, Proposition~\ref{prop:op-branchwise-refinement}, and Proposition~\ref{prop:op-admissible-routing}. Statements (iii)--(iv) are the corresponding compressed consequences of Theorem~\ref{thm:op3-shell} and Theorem~\ref{thm:combined-branch-shell}, with Proposition~\ref{prop:packet-phase-epsilon-screening} supplying the current diagnostic preselection reading of the packet/phase-lock/epsilon layer.
\end{proof}

\paragraph{Interpretation and reading guide.}
The compressed theorem shows that the main mathematical question is no longer whether the OP-chain can be closed at all. That closure is already fixed at the source-return level. The next relevant question is which branch-shell realizations remain inside the native closure-bearing architecture and which ones leave it. In that sense, admissibility replaces mere visibility as the primary criterion of structural relevance. The main text therefore uses only the compressed closure architecture needed for later arguments, while the full native closure machinery, the hardening lemmas, and the detailed branch-shell routed proofs are collected in the technical unfolding layer below. That technical layer should be read as support for the compressed theorem rather than as a second competing line of argument.

\subsection{OP3 as residual/shell interface}
\label{sec:op3-interface-main}
Residual and shell-sensitive structure is localized at the OP~3 interface. Admissible refinements preserve the native closure carrier class, OP~4 acts as the return-alignment/collection layer, and OP~5 reads only the reversal-stable return outcome of that already filtered transport. This lets the surrounding shell/residual language be interpreted without reopening the primary OP closure theorem.

\subsection{Combined branch-shell consequence}
\label{sec:combined-branch-shell-main}
A visibly composite return sector may be assembled from several closure-admissible branches and several OP~3 shell-admissible refinements without destroying structural OP closure. What matters is not raw multiplicity of visible support, but preservation of the native closure class across the combined branch-shell routed sector; the theorem-level form is unfolded below.

